\documentclass[11pt]{article}
\usepackage[a4paper, top=18mm, bottom=18mm, left=20mm, right=20mm]{geometry}
\usepackage[T2A]{fontenc}
\usepackage[utf8]{inputenc}
\usepackage[ukrainian]{babel}
\usepackage{graphicx}
\usepackage[export]{adjustbox}
\usepackage{amsmath}
\usepackage{amssymb}
\setlength{\parindent}{0pt}
\setlength{\parskip}{4pt}
\pagestyle{empty}
\begin{document}
%begin
{\large\textbf{Контрольна робота}}

\textbf{Варіант \No\,4}\hfill \textit{ПІБ:}\,\underline{\hspace{60mm}}
\vspace{4mm}

%beginitem
1. 
Значенням змінної \kw{a} є цілочисловий масив типу $numpy.ndarray$ розмірів $5{\times}5$. На скільки збільшиться сума елементів масиву \kw{a} після виконання
\begin{verbatim}
a.T[2] += 3
\end{verbatim}


Якщо код некоректний, то в якості відповіді зазначити ERROR.



%enditem

%beginitem
2. Що буде виведено за виконання фрагменту коду?

%code
cout << (false != 7 > 4 < 2);
%code

%enditem

%beginitem
3. Що буде виведено за виконання фрагменту коду?

%code
#include <iostream>
using namespace std;

int h(int a, int b){
    int c = 84;
    if (b != -3)
        c = 5;
    else if (a > -2)
         return 2;
    else 
        return 0;
    return c;
}

int main(){
    cout << h(3, 3);
    return 0;
}
%code

%enditem

%beginitem
4. Що буде виведено за виконання фрагменту коду?

%code
print('R', end=' ')
f=[0,1,2,3,4,5,6,7,8,9]
print(f[:-3:-3])
%code


%enditem

%beginitem
5. %goodpath:Scripts/2019/Mod2019_1/Lex/01-good.txt
Відмітити все, що є однією правильною лексемою:
( 1 ) 5_w

( 2 ) 0x17

( 3 ) 'c'

( 4 ) "1,3,2"

%enditem

%beginitem
6. Що буде виведено за виконання фрагменту коду?

%code
a,b,c=5,2,0
def h(b):
    a=3
    b=4
    c=2
    return a+b+c

a,b,c=9,6,8
print(h(b),a,b,c)
%code

%enditem

%beginitem
7. 
try:
    print(5, end=';')
    print(2j!=5j, end=';')
    print(0, end=';')
except TypeError:
    print(9, end=';')
except ZeroDivisionError:
    print(6, end=';')
except:
    print('Z', end=';')
else:
    print(8, end=';')
finally:
    print(1)

%enditem

%beginitem
8. 
Записати ВИРАЗ, значенням якого є множина, що складається з кубівцілих чисел від 12 до 2012 (включно), що не діляться на жодне з чисел 2, 3.
%enditem

%beginitem
9. %goodpath:Scripts/2018/Mod2018_1/01-good.txt
Відмітити все, що є однією правильною лексемою:
( 1 ) 5_w

( 2 ) (1)

( 3 ) )

( 4 ) 0.5

%enditem

%beginitem
10. %goodpath:Scripts/2019/Mod2019_1/Lex/03-good.txt
Відмітити все, що є коректним оператором С++:
( 1 ) x not is y

( 2 ) print(end="1","qwe")

( 3 ) 3**-1

( 4 ) x-=2

%enditem

%end
\newpage
\end{document}

